\section{Experience}

\ResumeSimpleEntry{HTC Viverse \textbar{} Senior Engineer @ Research Team}{May. 2022 -- July. 2025}
\begin{resumeitems}
\item \textbf{Global Platform Scale}: Served as the Lead Architect for the scalable backend design for the Viverse Worlds global launch. Conceptualized a dynamic URL resolution strategy, partnering with Ops on the edge infrastructure (AWS Lambda@Edge) to scale to 1 million monthly-active users.
\item \textbf{Algorithmic R\&D (Patented)}: Engineered and patented novel spatial computing algorithms and Machine Learning (Reinforcement Learning) frameworks to dynamically evaluate occlusion, successfully slashing distributed system network egress bandwidth by up to 80\%.
\item \textbf{AI Integration \& Revenue Growth}: Architected an event-driven AWS sidecar for low-latency WebRTC streams, enabling secure third-party AI moderation and driving a \$33,000 USD revenue increase.
\item \textbf{Cross-Functional Technical Leadership}: Directed multi-disciplinary engineering pods (Research, BE/FE, Ops) from system design to execution, transitioning the core platform from an isolated, single-scene architecture into an interconnected 3D ecosystem while maintaining strict backward compatibility.
\item \textbf{Concurrency Optimization}: Engineered a multi-worker paradigm in the WebRTC core to break single-thread CPU limits, scaling concurrent room capacity from 300 to 4,500+ consumers.
\end{resumeitems}

\ResumeSimpleEntry{HTC Vive \textbar{} Senior Engineer}{Sep. 2020 -- May. 2022}
\begin{resumeitems}
\item \textbf{Proactive Leadership \& Standardization}: Autonomously identified legacy I/O bottlenecks and drove a high-performance microservice evolution. Authored an independent, open-source FastAPI boilerplate that HTC Vive officially adopted as its foundational backend standard, reducing new API latency by 85\%.
\item \textbf{Asynchronous Architecture}: Architected a fault-tolerant Backend-For-Frontend (BFF) layer utilizing Celery and AWS SQS, fully decoupling CPU-heavy 3D asset manipulation to guarantee high availability and eliminate API Gateway timeouts.
\item \textbf{Multi-Tenant Security}: Engineered a highly scalable Role-Based Access Control (RBAC) ecosystem leveraging Domain-Driven Design (DDD) and polyglot persistence to securely isolate enterprise client data across a multi-tenant platform.
\end{resumeitems}

\ResumeSimpleEntry{InfoBoom \textbar{} Senior Software Developer (Technical Lead)}{Feb. 2020 -- Sep. 2020}
\begin{resumeitems}
\item \textbf{Technical Team Leadership}: Directed a 4-person engineering pod through complex architectural modernizations. Managed critical-path dependencies and unblocked technical constraints to guarantee the on-time delivery of high-security SaaS milestones.
\item \textbf{Cryptographic Architecture}: Engineered the high-risk modernization of a core C++ encryption library for a high-security SaaS platform, achieving strict zero-regression API compatibility while securing ISO27001 and Taiwan MAS compliance.
\item \textbf{System Modernization}: Architected the enterprise teardown of a legacy WebRTC monolith, migrating the infrastructure into a highly available, event-driven Node.js microservice ecosystem utilizing MongoDB and RabbitMQ.
\end{resumeitems}

\ResumeSimpleEntry{Synology \textbar{} Product Developer}{Dec. 2016 -- Jun. 2018}
\begin{resumeitems}
\item \textbf{Core Systems Engineering}: Engineered full-stack mail flow control features natively in C++ (Postfix/Dovecot), resolving critical system-level memory leaks and race conditions.
\end{resumeitems}
